% ============================================================
% MÓDULO: VIÁTICOS, RENDICIONES Y CAJA CHICA
% ============================================================
\clearpage
\section*{MÓDULO: VIÁTICOS, RENDICIONES Y CAJA CHICA}
\subsection*{El técnico en comisión de servicio}

\begin{recuadro}
\textbf{Objetivo del módulo:} saber cómo manejar el dinero de la empresa cuando se trabaja
fuera del taller: pedir y usar viáticos, solicitar los comprobantes correctos, armar una
rendición de gastos, administrar una caja chica y organizar un viaje de trabajo.
\end{recuadro}

% ------------------------------------------------------------
\section*{1. ¿QUÉ SON LOS VIÁTICOS?}

Un técnico aviónico muchas veces trabaja \textbf{fuera de su lugar habitual}: un AOG
(\textit{Aircraft On Ground}, avión detenido en tierra por una falla) en otro aeropuerto, una
instalación en un aeroclub del interior, una inspección en un hangar de un cliente. Los gastos
de ese trabajo (pasajes, combustible, peajes, comida, alojamiento, materiales comprados en el
lugar) los paga la empresa. Esos gastos se llaman \textbf{viáticos}.

Hay dos modalidades:
\begin{itemize}
  \item \textbf{Anticipo de viáticos:} la empresa entrega dinero (o una tarjeta) antes del
        viaje, y el técnico después \textbf{rinde} lo que gastó y devuelve lo que sobró.
  \item \textbf{Reintegro:} el técnico paga con su dinero y después la empresa le devuelve lo
        gastado \textbf{contra presentación de comprobantes}.
\end{itemize}

\paragraph*{Viáticos y sueldo.}
Los viáticos \textbf{no son sueldo}: son la devolución de un gasto hecho por cuenta de la
empresa. Por eso, cuando están respaldados con comprobantes, no se consideran remuneración (no
llevan aportes). En cambio, si la empresa paga ``viáticos'' fijos sin comprobantes, en la
práctica pueden ser considerados parte del salario.

% ------------------------------------------------------------
\section*{2. CÓMO PEDIR LOS COMPROBANTES}

Para que la empresa pueda justificar el gasto (y, si es Responsable Inscripta, descontar el
IVA y deducirlo en Ganancias), el comprobante tiene que estar \textbf{a nombre de la empresa},
no del técnico.

\begin{enumerate}
  \item Antes de viajar, tener a mano los \textbf{datos de facturación} de la empresa: razón
        social, CUIT, condición frente al IVA y domicilio. Conviene guardarlos en el celular.
  \item Al pagar, pedir: \textit{``¿Me hacés factura A a nombre de la empresa? Te paso el
        CUIT''}.
  \item Qué factura corresponde:
  \begin{itemize}
    \item Si la empresa es \textbf{Responsable Inscripta} y el comercio también:
          \textbf{factura A} (con el IVA discriminado).
    \item Si el comercio es \textbf{monotributista}: \textbf{factura C} a nombre de la
          empresa.
    \item Si la empresa es monotributista: \textbf{factura B} (si el comercio es RI) o C.
  \end{itemize}
  \item \textbf{Ticket de controlador fiscal:} en estaciones de servicio y comercios, se puede
        pedir \textbf{ticket factura A} dando el CUIT antes de que emitan el ticket.
  \item Revisar el comprobante \textbf{en el momento}: CUIT correcto, fecha, importe y que
        tenga CAE (código de autorización de ARCA) si es electrónico.
\end{enumerate}

\begin{recuadro}
\textbf{Gastos sin factura.} Algunos gastos no tienen factura: peajes con telepeaje (se
descarga el resumen), estacionamiento callejero, propinas, transporte informal. Se registran
en una \textbf{planilla de gastos sin comprobante} con fecha, concepto, importe y motivo, y la
empresa decide si los acepta. Siempre son la excepción: \textbf{sin comprobante, el gasto
puede no reintegrarse}.
\end{recuadro}

% ------------------------------------------------------------
\section*{3. MEDIOS DE PAGO DE LA EMPRESA}

\begin{itemize}
  \item \textbf{Efectivo:} el anticipo clásico. Exige guardar todos los comprobantes y devolver
        el sobrante.
  \item \textbf{Tarjetas prepagas o recargables corporativas:} la empresa carga un monto en
        una tarjeta a nombre del técnico. Todos los gastos quedan registrados, se fijan
        límites y se pueden bloquear rubros. Hay que rendir igual: el movimiento de la tarjeta
        no reemplaza a la factura.
  \item \textbf{Tarjeta de crédito corporativa:} para técnicos o gerentes que viajan seguido;
        la deuda es de la empresa.
  \item \textbf{Transferencia a la cuenta del técnico:} para anticipos o reintegros.
\end{itemize}
\textbf{Regla de oro: nunca mezclar el dinero de la empresa con el propio.} No usar la tarjeta
de la empresa para gastos personales, ni siquiera ``para devolverlo después''.

% ------------------------------------------------------------
\section*{4. LA RENDICIÓN DE GASTOS}

La \textbf{rendición} es el informe en el que el técnico detalla qué gastó, en qué y con qué
comprobantes. Hoy se hace con planillas o con aplicaciones de la empresa (sacando foto a cada
comprobante), muchas veces integradas al ERP.

\subsubsection*{4.1. Modelo de planilla de rendición}
\begin{center}
\footnotesize
\setlength{\tabcolsep}{3pt}
\begin{tabular}{|c|p{0.10\linewidth}|p{0.17\linewidth}|p{0.14\linewidth}|p{0.11\linewidth}|p{0.15\linewidth}|r|}
\hline
\textbf{N°} & \textbf{Fecha} & \textbf{Concepto} & \textbf{Proveedor} & \textbf{Comprob.}
& \textbf{Medio de pago} & \textbf{Importe} \\ \hline
1 & 03/06 & Combustible & YPF Ruta 2 & Fact. A & Tarj. prepaga & \$45.000 \\ \hline
2 & 03/06 & Peajes & AUBASA & Resumen & Telepeaje & \$6.800 \\ \hline
3 & 03/06 & Almuerzo & Parador & Fact. C & Efectivo & \$14.500 \\ \hline
4 & 03/06 & Conector BNC (2 u.) & Electrónica Sur & Fact. A & Tarj. prepaga & \$9.200 \\
\hline
5 & 03/06 & Hotel 1 noche & Hotel Centro & Fact. A & Tarj. prepaga & \$62.000 \\ \hline
\multicolumn{6}{|r|}{\textbf{Total gastado}} & \textbf{\$137.500} \\ \hline
\multicolumn{6}{|r|}{Anticipo recibido} & \$150.000 \\ \hline
\multicolumn{6}{|r|}{\textbf{Saldo a devolver a la empresa}} & \textbf{\$12.500} \\ \hline
\end{tabular}
\end{center}

\subsubsection*{4.2. Buenas prácticas}
\begin{itemize}
  \item Indicar el \textbf{motivo del viaje} y la \textbf{orden de trabajo} o el cliente al
        que se imputa el gasto (en un taller, los gastos se pueden cobrar al cliente).
  \item Numerar los comprobantes y pegarlos o escanearlos en el mismo orden de la planilla.
  \item Rendir \textbf{en el plazo} que fija la empresa (por ejemplo, dentro de los 5 días
        hábiles del regreso).
  \item Respetar la \textbf{política de viáticos}: topes diarios de comida, categoría de hotel,
        qué gastos se aceptan y cuáles no (por ejemplo, bebidas alcohólicas).
  \item Los papeles térmicos se borran: sacar foto al comprobante el mismo día.
\end{itemize}

% ------------------------------------------------------------
\section*{5. CAJA CHICA}

La \textbf{caja chica} (o fondo fijo) es una pequeña suma de dinero que se deja en el taller o
en el pañol, a cargo de un \textbf{responsable}, para pagar gastos menores y urgentes: un
componente que falta, insumos de limpieza, un envío, un taxi.

\paragraph*{Sistema de fondo fijo.}
\begin{enumerate}
  \item La empresa asigna un monto fijo (por ejemplo, \$200.000) y un responsable.
  \item Cada gasto se paga con la caja y se guarda el comprobante (a nombre de la empresa).
  \item En todo momento: \textbf{efectivo en caja + comprobantes = monto del fondo}.
  \item Cuando el efectivo baja, el responsable \textbf{rinde} los comprobantes y la empresa
        \textbf{repone} exactamente lo gastado, volviendo al monto fijo.
  \item Se hacen \textbf{arqueos} (recuentos) sorpresivos para controlar.
\end{enumerate}

% ------------------------------------------------------------
\section*{6. ORGANIZAR UN VIAJE DE TRABAJO}

\begin{itemize}
  \item \textbf{Antes de salir:} confirmar la orden de trabajo, dirección exacta y contacto en
        destino; preparar herramientas, instrumental y repuestos (lista de verificación); pedir
        el anticipo o la tarjeta; llevar DNI, credencial o licencia y, si se ingresa a
        plataforma, el permiso de acceso aeroportuario.
  \item \textbf{Transporte:} auto de la empresa (con seguro, cédula y documentación al día),
        micro de larga distancia, avión o remís. Siempre pedir comprobante a nombre de la
        empresa. Para herramientas en avión: consultar las restricciones de equipaje.
  \item \textbf{Alojamiento:} reservar con factura a nombre de la empresa; elegir cerca del
        lugar de trabajo.
  \item \textbf{Durante el trabajo:} registrar horas trabajadas, horas de viaje y lo realizado
        (informe técnico para el cliente).
  \item \textbf{Al volver:} devolver herramientas al pañol, rendir los gastos y entregar el
        informe.
\end{itemize}

% ------------------------------------------------------------
\section*{7. ACTIVIDADES DEL MÓDULO}

\begin{enumerate}
  \item \textbf{Simulación de AOG:} un avión quedó en tierra en un aeropuerto a 400 km por una
        falla de un equipo de comunicaciones. Planificar el viaje (transporte, alojamiento,
        herramientas y repuestos) y estimar el anticipo necesario.
  \item \textbf{Rendición:} con un sobre de comprobantes de ejemplo (facturas A, B y C,
        tickets, un resumen de telepeaje, un gasto sin comprobante), armar la planilla de
        rendición y calcular el saldo a devolver o a reintegrar. Identificar qué comprobantes
        tienen errores.
  \item \textbf{Pedir la factura:} representar en parejas cómo pedir una factura A a nombre de
        la empresa en un comercio, en una estación de servicio y en un hotel.
  \item \textbf{Caja chica del pañol:} definir el monto, el responsable y las reglas de uso de
        una caja chica para el pañol de la escuela, y registrar sus movimientos durante un mes
        (puede integrarse al proyecto Odoo).
\end{enumerate}
